\chapter{Tag 9 — Wir zeichnen einen Strichcode mit der Hand}

\begin{quote}
Bisher hat Python jede Aufgabe abgenommen, sobald sie mathematisch
verstanden war. Heute drehen wir den Spieß um: du nimmst Bleistift,
Lineal und einen schwarzen Filzstift und zeichnest einen
EAN-13-Strichcode \emph{Modul für Modul} auf Papier. Am Ende soll
dein Handy ihn abscannen — kein Drucker, kein Software-Helfer,
nur dein Wissen aus den letzten Tagen plus eine ruhige Hand.
\end{quote}

\section*{Lernziele für heute}

Am Ende von Tag 9 kannst du \dots

\begin{itemize}
\item die drei Encoding-Sets eines EAN-13 lesen: L (links gerade), G
  (links ungerade), R (rechts)
\item erklären, wie die führende Ziffer das Paritätsmuster der linken
  Hälfte steuert — und warum dieser Umweg überhaupt nötig ist
\item einen 13-stelligen EAN auf eine vorbereitete Zeichenvorlage
  übertragen — Modul für Modul, Strich für Strich
\item den selbstgezeichneten Strichcode mit dem Handy scannen
\item erklären, warum EAN-13 \emph{nur} Fehler erkennen kann, aber
  keine korrigieren — und warum Datamatrix das anders macht
\end{itemize}

\paragraph{Material} Ein Lineal mit Millimeter-Skala, ein Fineliner
(Strichbreite 0{,}3 oder 0{,}4\,mm) und die Zeichenvorlage zu dieser
Etappe — ein separates A4-Heft mit \emph{mehreren} Barcode-Feldern.
Mehrere Felder deshalb, weil der erste Versuch fast nie auf Anhieb
sitzt; falls doch, hast du Reserve für eigene Codes.

% =============================================================
\section{Prüfziffer nachrechnen am eigenen EAN \normalfont(\textit{≈ 15 min})}
% =============================================================

An Tag 1 hast du die EAN-13-Prüfziffer kennengelernt. Heute brauchst
du wieder einen konkreten 13-stelligen Code — denn diesen Code wirst
du gleich \emph{auf Papier zeichnen}, und dafür muss er stimmen.

Wir nehmen für die ganze Etappe einen ganz bestimmten Code: den von
einem 500-Gramm-Glas Honig „Sommer in Radebeul“ von \textbf{Gommels
Bienen}. Auf der Rückseite des Glases steht der EAN

\begin{center}
{\Large\ttfamily 4\,270004\,371635}
\end{center}

\begin{bleistiftuebung}{1}{Prüfziffer nachrechnen}
Verifiziere die Prüfziffer der EAN \texttt{4270004371635} nach der
Methode aus Tag 1: Decke die letzte Ziffer ab, berechne sie aus den
ersten zwölf Ziffern, und vergleiche mit der gedruckten Ziffer.

Im Lösungsanhang findest du die komplette Schritt-für-Schritt-Rechnung —
\textbf{bitte erst dort nachsehen}, nachdem du es selbst probiert hast.
\end{bleistiftuebung}

Wenn das Ergebnis deiner Methode mit der gedruckten Prüfziffer
übereinstimmt, weißt du: der Code, den du gleich zeichnen wirst, ist
\emph{technisch sauber}. Nichts ist frustrierender, als am Ende einer
Stunde Zeichnen einen Strichcode in den Händen zu halten, den dein
Handy abweist, weil eine Ziffer falsch war.

% =============================================================
\section{Drei Codetabellen — und warum es drei sein müssen \normalfont(\textit{≈ 25 min})}
% =============================================================

So sieht das Ergebnis aus, das du am Ende dieser Etappe auf Papier
gezeichnet haben sollst:

\begin{center}
\barcodehonig
\end{center}

Das ist der EAN-13 unseres Honigglases, gerendert mit derselben Logik,
die du gleich kennenlernst. Weiße Module = leerlassen, schwarze Module =
schwarz machen. So einfach — und so geduldig.

\subsection*{Die kleine Geschichte des EAN-13}

In den USA wurde 1973 der \textbf{UPC-A} eingeführt — ein 12-stelliger
Strichcode für Konsumgüter. Layout: 6 Ziffern links, 6 Ziffern rechts,
zwei Codetabellen (eine für jede Hälfte). Funktionierte gut.

Drei Jahre später wollte Europa mitspielen, brauchte aber eine
\emph{13.} Ziffer für die Länderkennung (4 = Deutschland, 8 =
Italien, etc.). Es gab nur ein Problem: Der UPC-A-Scanner-Park war
schon ausgerollt. Niemand wollte alle Scanner umrüsten.

Die Lösung der EAN-13-Erfinder ist eine elegante Pirouette: Sie haben
die Strich-Geometrie \emph{nicht} verändert. Es bleibt bei 12
Strich-Positionen. Stattdessen führten sie auf der linken Seite eine
\emph{zweite} Codetabelle ein. Pro Ziffer kann man jetzt zwischen
zwei Mustern wählen — L oder G — und welches Muster wo verwendet
wird, codiert die 13. Ziffer „unsichtbar" mit.

Damit das alte UPC-A weiterhin funktioniert, wurde die Konvention
gesetzt: Ein EAN-13, der mit \texttt{0} beginnt, hat alle linken
Ziffern in L-Codierung — und sieht für einen UPC-A-Scanner exakt aus
wie ein UPC-A.

\subsection*{Das Modul als Bauteil}

Bevor wir uns die Tabellen anschauen: Ein EAN-13 besteht aus
\textbf{Modulen} — gleich breiten schwarzen oder weißen Streifen. Die
Modulbreite ist genormt; bei einem klassischen Strichcode auf einer
Konsumverpackung sind das etwa $0{,}33$\,mm. Ein schwarzer Strich kann
ein, zwei, drei oder vier Module breit sein — entscheidend ist nur,
wie viele Module nebeneinander schwarz sind, und wie viele weiß.

Jede Ziffer wird durch \textbf{7 Module} dargestellt — eine Folge
aus 7 Bits, wobei \texttt{0} für weiß und \texttt{1} für schwarz
steht. Welche der drei Tabellen man verwendet (L, G oder R), hängt
davon ab, an welcher Position die Ziffer im Strichcode steht.

\subsection*{Die L-Codes (linke Hälfte, ungerade Parität)}

\emph{L} steht für \emph{left-hand odd-parity}: Codes für die linke
Hälfte mit einer ungeraden Anzahl schwarzer Module. Jeder L-Code
beginnt mit \texttt{0} (weiß) und endet mit \texttt{1} (schwarz).

\begin{center}
\begin{tabular}{c c c c c}
\toprule
\textbf{Ziffer} & \textbf{L-Code} & \textbf{G-Code} & \textbf{R-Code} & \textbf{schwarze Module}\\
\midrule
0 & \texttt{0001101} & \texttt{0100111} & \texttt{1110010} & 3 / 4 / 4 \\
1 & \texttt{0011001} & \texttt{0110011} & \texttt{1100110} & 3 / 4 / 4 \\
2 & \texttt{0010011} & \texttt{0011011} & \texttt{1101100} & 3 / 4 / 4 \\
3 & \texttt{0111101} & \texttt{0100001} & \texttt{1000010} & 5 / 2 / 2 \\
4 & \texttt{0100011} & \texttt{0011101} & \texttt{1011100} & 3 / 4 / 4 \\
5 & \texttt{0110001} & \texttt{0111001} & \texttt{1001110} & 3 / 4 / 4 \\
6 & \texttt{0101111} & \texttt{0000101} & \texttt{1010000} & 5 / 2 / 2 \\
7 & \texttt{0111011} & \texttt{0010001} & \texttt{1000100} & 5 / 2 / 2 \\
8 & \texttt{0110111} & \texttt{0001001} & \texttt{1001000} & 5 / 2 / 2 \\
9 & \texttt{0001011} & \texttt{0010111} & \texttt{1110100} & 3 / 4 / 4 \\
\bottomrule
\end{tabular}
\end{center}

\paragraph{Beziehung zwischen den drei Tabellen}
\begin{itemize}
\item \textbf{R-Code} = L-Code, jedes Bit invertiert (aus \texttt{0}
  wird \texttt{1} und umgekehrt). Schwarz und Weiß tauschen.
\item \textbf{G-Code} = R-Code, rückwärts gelesen. Spiegelung.
\item Daraus folgt: G hat dieselbe Bit-Anzahl wie R, aber die L- und
  G-Codes lassen sich \emph{nicht} ineinander umwandeln, ohne über R
  zu gehen.
\end{itemize}

Du musst die Tabelle nicht auswendig können — sie steht in der
Zeichenvorlage unten auf jedem Blatt.

\subsection*{Die führende Ziffer als Paritätsmuster}

Auf der linken Hälfte des EAN-13 stehen die Ziffern 2 bis 7 (sechs
Stück). Pro Ziffer wird entweder L oder G verwendet — und genau diese
Wahl, sechs Mal hintereinander, codiert die führende (erste)
Ziffer:

\begin{center}
\begin{tabular}{c l}
\toprule
\textbf{Erstziffer} & \textbf{L/G-Muster der nächsten 6 Ziffern} \\
\midrule
0 & \texttt{LLLLLL} \\
1 & \texttt{LLGLGG} \\
2 & \texttt{LLGGLG} \\
3 & \texttt{LLGGGL} \\
4 & \texttt{LGLLGG} \\
5 & \texttt{LGGLLG} \\
6 & \texttt{LGGGLL} \\
7 & \texttt{LGLGLG} \\
8 & \texttt{LGLGGL} \\
9 & \texttt{LGGLGL} \\
\bottomrule
\end{tabular}
\end{center}

Auf der rechten Hälfte stehen die Ziffern 8 bis 13 (auch sechs
Stück) — die werden immer in R-Codierung gesetzt, ohne Variation.

\subsection*{Beispiel: unser Honig-EAN}

Unser Code ist \texttt{4270004371635}. Die führende Ziffer ist
\textbf{4}, also lautet das Paritätsmuster \texttt{LGLLGG}. Daraus
folgt für die ersten sieben sichtbaren Ziffern (Position 2 bis 7) die
Codetabelle:

\begin{center}
\begin{tabular}{c c c c c c c}
\toprule
\textbf{Position} & 2 & 3 & 4 & 5 & 6 & 7 \\
\textbf{Ziffer}   & 2 & 7 & 0 & 0 & 0 & 4 \\
\textbf{Tabelle}  & L & G & L & L & G & G \\
\bottomrule
\end{tabular}
\end{center}

Die rechten Ziffern \texttt{371635} werden alle in R-Codierung
gesetzt. Die führende Ziffer \texttt{4} wird \emph{nicht} als Strich
gezeichnet — sie ist nur unten in Klarschrift sichtbar, ihre
Information liegt im L/G-Muster der linken Hälfte.

\begin{bleistiftuebung}{2}{Encoding-Folge ablesen}
Welche Codetabellen-Folge ergibt sich für einen EAN, der mit
\texttt{9} beginnt? Und für einen, der mit \texttt{0} beginnt — und
warum sieht der dann „wie ein UPC-A" aus?
\end{bleistiftuebung}
